\documentclass[11pt]{article}
\usepackage{preamble}
\setlength{\parskip}{4pt}

\begin{document}

\daybanner{AP Statistics \textperiodcentered\ Unit 5 \textperiodcentered\ Topic 5.3 \textperiodcentered\ B: 2027-02-09 \textperiodcentered\ E: 2027-04-05 \textperiodcentered\ TEACHER KEY}{Two Payloads, One Regression Line}

\sectionbanner{CED TETHER}

{\small
\begin{itemize}[leftmargin=1.5em,itemsep=2pt,topsep=2pt]
  \item \textbf{Skill 2.C} --- Calculate summary statistics, relative positions of points within a distribution, correlation, and predicted response.
  \item \textbf{EU DAT-1} --- Regression models may allow us to predict responses to changes in an explanatory variable.
  \item \textbf{LO DAT-1.D} --- Calculate a predicted response value using a linear regression model. [Skill 2.C]
  \item \textbf{EK DAT-1.D.1} --- A simple linear regression model is an equation that uses an explanatory variable, x, to predict the response variable, y.
  \item \textbf{EK DAT-1.D.2} --- The predicted response value, denoted by $\widehat{y}$, is calculated as $\widehat{y}$ = a + bx, where a is the y-intercept and b is the slope of the regression line, and x is the value of the explanatory variable.
  \item \textbf{EK DAT-1.D.3} --- Extrapolation is predicting a response value using a value for the explanatory variable that is beyond the interval of x-values used to determine the regression line. The predicted value is less reliable as an estimate the further we extrapolate.
\end{itemize}
}
\textbf{Essential question:} How should observed range, prediction purpose, and new evidence guide use of a regression model?\par
\textbf{DOK-3 skill:} 4.B \textperiodcentered\ \textbf{FRQ pattern:} {\small\texttt{regression-\allowbreak{}prediction-\allowbreak{}range-\allowbreak{}and-\allowbreak{}decision}}\par
\textbf{DOK rationale:} Part (c) weighs competing claims using the day's numerical or graphical evidence and explains a limitation or consequence; the judgment requires linked reasoning beyond a calculation.\par

\frameworkphaseheader{Do Now (first take) $\rightarrow$ Explore (video + follow-along) $\rightarrow$ Exit (finish + turn in)}{1 $\rightarrow$ 3}{45}{%
\begin{itemize}[leftmargin=1.5em,itemsep=2pt,topsep=2pt]
  \item Show the board at the bell and collect a one-sentence first take without confirming a claim.
  \item Run the named video follow-along, then allow about 10 minutes for parts (a)--(c).
  \item Ask for evidence linked to a claim. Collect the sheet and hand-score part (c) with E/P/I.
\end{itemize}
}{%
\begin{itemize}[leftmargin=1.5em,itemsep=2pt,topsep=2pt]
  \item Commit to a recommendation and cite one number or plot feature.
  \item Complete the follow-along, finish (a)--(c), revise, and turn in.
\end{itemize}
}{%
\begin{itemize}[leftmargin=1.5em,itemsep=2pt,topsep=2pt]
  \item Which number or visible feature supports your claim?
  \item Why does the other claim go beyond that evidence?
  \item What would you tell someone who chose the other claim?
\end{itemize}
}{Circulate and ask students to connect their choice to evidence. Score part (c) using the three required reasoning elements; do not require dropped extension work.}

\begin{callout}{\IconWarn\ WATCH FOR}{calloutred}
\begin{itemize}[leftmargin=1.5em,itemsep=2pt,topsep=2pt]
  \item A choice with copied numbers but no explanation of how they support it.
  \item Treating a model or average as a guaranteed individual outcome.
\end{itemize}
\end{callout}

\sectionbanner{SCORING PART (c) --- E / P / I}

\textbf{Expected elements}
\begin{itemize}[leftmargin=1.5em,itemsep=2pt,topsep=2pt]
  \item \textbf{reasoning-1} (required) --- Makes a supported decision by comparing an in-range and out-of-range prediction
  \item \textbf{reasoning-2} (required) --- Uses the tested range and 1.60-minute predicted surplus without treating it as guaranteed
  \item \textbf{reasoning-3} (required) --- Requests flight evidence near 3.5 kg before relying on the model
\end{itemize}

{\small\renewcommand{\arraystretch}{1.25}
\noindent\begin{tabular}{|p{0.5in}|p{5.9in}|}
\hline
\textbf{E} & Includes all three required reasoning elements above, with correct contextual evidence. \\ \hline
\textbf{P} & Includes at least one substantive correct reasoning link above, but omits or misstates another required element. \\ \hline
\textbf{I} & Includes no substantive correct reasoning link above; an unsupported choice or copied number alone is insufficient. \\ \hline
\end{tabular}}

\clearpage
\focusbanner{TODAY'S PROBLEM --- annotated}

Suppose eight drone tests with payloads from 0.5 to 2.5 kg produced the shown data and the model $\hat y=36.0-6.4x$, where $x$ is payload and $\hat y$ is predicted flight duration in minutes. A proposed mission carries 3.5 kg and requires 12 minutes of flight.\par\smallskip \textbf{Ravi:} Approve it: the model predicts more than the required 12 minutes, and the fitted data are close to a line.\par \textbf{Mina:} Before deciding, compare the support for predictions at 1.8 and 3.5 kg and identify evidence the safety decision still needs.\par
\begin{center}
\begin{tikzpicture}
\begin{axis}[
  width=0.52\linewidth,
  height=1.40in,
  xmin=0.3, xmax=2.7, ymin=18, ymax=35,
  xlabel={Payload (kg)}, ylabel={Duration (min)}, grid=major,
  tick label style={font=\small}, label style={font=\small}
]
\addplot[only marks, mark=*, mark size=1.3pt, color=dayblue] coordinates {
  (0.5,33.2)
  (0.8,30.58)
  (1.1,29.16)
  (1.4,26.74)
  (1.6,25.46)
  (1.9,24.04)
  (2.2,21.62)
  (2.5,20.4)
};
\end{axis}
\end{tikzpicture}
\end{center}
\begin{firsttakebox}{FIRST TAKE (5 min)}
Would you approve now, delay approval, or approve conditionally? Cite one number or plot feature.\par
\answer{Not graded. Students should distinguish a computed prediction from evidence supporting its use.}
\end{firsttakebox}

\textit{Rules callout ``USING A REGRESSION MODEL (keep on the page)'' is printed on the student sheet.}\par\medskip

\dokbadge{1}~\textbf{(a)} Name the explanatory and response variables, with units.\par
\answer{Payload (kg) is explanatory; flight duration (minutes) is the response.}

\dokbadge{2}~\textbf{(b)} Predict flight duration at 1.8 and 3.5 kg. How much does the 3.5-kg prediction exceed the required 12 minutes?\par
\answer{At 1.8 kg, $\hat y=36-6.4(1.8)=24.48$ minutes. At 3.5 kg, $\hat y=36-6.4(3.5)=13.60$ minutes, leaving $13.60-12=1.60$ minutes above the requirement.}

\dokbadge{3}~\textbf{(c)} Would you approve this mission now? Compare the support for the two predictions using the tested payload range and the extra time from (b). Explain what evidence you would want before relying on the 3.5-kg prediction. Write 3--4 sentences.\par
\answer{Mina's approach supports delaying approval. The 1.8-kg prediction is within the tested 0.5--2.5 kg range, but 3.5 kg is 1.0 kg beyond it, where the pattern may change. A predicted 1.60-minute surplus does not guarantee enough time, even though the tested points follow a line. Repeated tests near 3.5 kg are needed to check actual flight times before approval.}

\begin{summaryexitbox}{EXIT}
One thing the video changed about my first take: \hrulefill
\end{summaryexitbox}

\end{document}
